\begin{derivation}[从\eqnrefs{eq:si-scattering-fourmomentum-r9}到\eqnrefs{eq:positive-energy-mass-shell-measure-r68}]
先从四维积分开始，并逐项检查其 Lorentz 变换性质：

\begin{equation*}
 I[F]\equiv
 \int\dd^4 p\,
 \delta(p^2-m^2c^2)\,\theta(p^0)F(p),
 \end{equation*}
其中 $\dd^4 p$ 的雅可比行列式在保向正时 Lorentz 变换下为 $1$，$p^2$ 是 Lorentz 标量，而这类变换保持 $p^0>0$。因此 $I[F]$ 对标量 $F$ 是 Lorentz 不变量。固定三动量后，把 $\delta$ 函数的自变量写成
\begin{equation*}
 f(p^0)=(p^0)^2-\mathbf p^2-m^2c^2.
\end{equation*}
两个根为 $p^0_\pm=\pm E_{\mathbf p}/c$，且
\begin{equation*}
 \left|\frac{\dd f}{\dd p^0}\right|_{p^0=p^0_\pm}
 =2\left|p^0_\pm\right|
 =\frac{2E_{\mathbf p}}{c}.
\end{equation*}
由一维恒等式
\begin{equation*}
 \delta[f(x)]
 =\sum_i\frac{\delta(x-x_i)}{|f'(x_i)|}
\end{equation*}
得到
\begin{equation*}
 \delta(p^2-m^2c^2)\theta(p^0)
 =\frac{c}{2E_{\mathbf p}}
 \delta\!\left(p^0-\frac{E_{\mathbf p}}{c}\right).
\end{equation*}
对 $p^0$ 积分以后，
\begin{equation*}
 I[F]
 =\int\frac{\dd^3 p}{2p^0}F(p)
 =\int\frac{c\,\dd^3 p}{2E_{\mathbf p}}F(p).
 \end{equation*}
左端是 Lorentz 不变量，因此右端就是正能质量壳上的不变量测度；完备关系~\eqnrefs{eq:rel-completeness-si-r9} 再按 Fourier 约定除以 $(2\pi\hbar)^3$。
\end{derivation}

\begin{derivation}[从\eqnrefs{eq:Ecm-s-si-r9}到\eqnrefs{eq:cm-momenta-kallen-si-r9}]
正文\eqnrefs{eq:Ecm-s-si-r9} 把质心总能量写成 $c\sqrt s$。在两体质心系中，两条粒子的三动量必定大小相同、方向相反，因此只剩一个未知的动量模。把两条质量壳条件与总能量守恒联立即可消去单粒子能量，得到正文\eqnrefs{eq:cm-momenta-kallen-si-r9}；根号何时变为实数同时给出二体阈值。

定义
\begin{equation*}
 \Lambda_K(a,b,c)=a^2+b^2+c^2-2ab-2ac-2bc.
 \end{equation*}
在初态质心系写
\begin{equation*}
 p_1^\mu=\left(\frac{E_1}{c},\mathbf p\right),
 \qquad
 p_2^\mu=\left(\frac{E_2}{c},-\mathbf p\right),
 \qquad
 c\sqrt s=E_1+E_2.
\end{equation*}
由
\begin{equation*}
 s=(p_1+p_2)^2=m_1^2c^2+m_2^2c^2+2p_1\!\cdot p_2
\end{equation*}
先得到
\begin{equation*}
 p_1\!\cdot p_2
 =\frac{s-(m_1^2+m_2^2)c^2}{2}.
\end{equation*}
另一方面
\begin{equation*}
 p_1\!\cdot p_2=\frac{E_1E_2}{c^2}+\mathbf p^2.
\end{equation*}
更直接的做法是先由 $P=p_1+p_2$ 与 $P^2=s$ 计算
\begin{align*}
 P\!\cdot p_1
 &=p_1^2+p_1\!\cdot p_2
 =\frac{s+(m_1^2-m_2^2)c^2}{2},\\
 P\!\cdot p_1
 &=\sqrt s\,\frac{E_1}{c},
\end{align*}
所以
\begin{equation*}
 E_1
 =\frac{c}{2\sqrt s}
 \left[s+(m_1^2-m_2^2)c^2\right].
\end{equation*}
再用 $E_1^2=m_1^2c^4+c^2p^2$，
\begin{align*}
 p^2
 &=\frac{1}{4s}
 \left[s+(m_1^2-m_2^2)c^2\right]^2-m_1^2c^2,\\
 &=\frac{\Lambda_K(s,m_1^2c^2,m_2^2c^2)}{4s}.
\end{align*}
因此
\begin{equation*}
 {
 |\mathbf p_i|
 =\frac{\sqrt{\Lambda_K(s,m_1^2c^2,m_2^2c^2)}}{2\sqrt s},
 \qquad
 |\mathbf p_f|
 =\frac{\sqrt{\Lambda_K(s,m_3^2c^2,m_4^2c^2)}}{2\sqrt s}.}
 \end{equation*}
末态阈值条件由根号非负给出 $c\sqrt s\ge(m_3+m_4)c^2$。
\end{derivation}

\begin{derivation}[从\eqnrefs{eq:phase2-si-start-r9}到\eqnrefs{eq:phase2-final-si-r9}]
正文\eqnrefs{eq:phase2-si-start-r9} 仍含两个末态三动量积分和一个四维守恒 $\delta$ 函数。在质心系中，三维 $\delta$ 函数先固定 $\mathbf p_4=-\mathbf p_3$，剩下的能量 $\delta$ 函数再固定径向动量模；最后只留下末态方向 $\dd\Omega$，从而得到正文\eqnrefs{eq:phase2-final-si-r9}。

采用能量四矢量
\begin{equation*}
 \mathcal P^\mu\equiv cp^\mu=(E,c\mathbf p),
 \end{equation*}
它只是物理四动量乘以常数 $c$。定义标准二体相空间因子
\begin{equation*}
 \dd\mathfrak P_2
 =(2\pi)^4\delta^{(4)}(\mathcal P_1+\mathcal P_2-\mathcal P_3-\mathcal P_4)
 \prod_{a=3,4}\frac{c^3\dd^3 p_a}{(2\pi)^3\,2E_a}.
 \end{equation*}
每个 $c^3\dd^3 p$ 都是 $\dd^3(c\mathbf p)$，所以 \eqnrefs{eq:phase2-si-start-r9} 完全由能量量纲的四动量组成。在质心系 $\mathcal P_1+\mathcal P_2=(c\sqrt s,\mathbf0)$。三维 $\delta$ 先令 $\mathbf p_4=-\mathbf p_3$，于是
\begin{equation*}
 \dd\mathfrak P_2
 =\frac{c^4p^2\dd p\,\dd\Omega}{(2\pi)^2 4E_3E_4}
 \delta(c\sqrt s-E_3-E_4).
\end{equation*}
令
\begin{equation*}
 f(p)=c\sqrt s-E_3(p)-E_4(p).
\end{equation*}
由 $\dd E_a/\dd p=c^2p/E_a$，
\begin{equation*}
 \left|f'(p)\right|
 =c^2p\left(\frac1{E_3}+\frac1{E_4}\right)
 =\frac{c^3p\sqrt s}{E_3E_4}.
\end{equation*}
所以
\begin{equation*}
 \delta[f(p)]
 =\frac{E_3E_4}{c^3p_f\sqrt s}\delta(p-p_f).
\end{equation*}
代回并积分 $p$，得到
\begin{equation*}
 {
 \dd\mathfrak P_2
 =\frac1{16\pi^2}\frac{|\mathbf p_f|}{\sqrt s}\,\dd\Omega.}
 \end{equation*}
右端无量纲，因为 $|\mathbf p_f|$ 与 $\sqrt s$ 都具有动量量纲。
\end{derivation}

\begin{derivation}[从\eqnrefs{eq:invariant-flux-scalar-r9}到\eqnrefs{eq:moller-velocity-si-r9}]
对两条入射轨迹，普通同一参考系中的相对相遇速率不能简单写成任意参考系都适用的 $|\mathbf v_1-\mathbf v_2|$。由四动量构造 Lorentz 标量

\begin{equation*}
 \mathcal I_{12}
 \equiv\sqrt{(p_1\!\cdot p_2)^2-m_1^2m_2^2c^4}.
 \end{equation*}
它具有动量平方量纲。定义 M\o ller 相对速率
\begin{equation*}
 v_{\rm M}
 \equiv
 \frac{c^3\mathcal I_{12}}{E_1E_2}.
 \end{equation*}
检查质心系时不跳过中间平方。写
\begin{equation*}
 p_1^\mu=\left(\frac{E_1}{c},\mathbf p_i\right),
 \qquad
 p_2^\mu=\left(\frac{E_2}{c},-\mathbf p_i\right),
\end{equation*}
于是
\begin{equation*}
 p_1\!\cdot p_2=\frac{E_1E_2}{c^2}+|\mathbf p_i|^2.
\end{equation*}
两条质量壳又分别给
\begin{equation*}
 m_a^2c^2=\frac{E_a^2}{c^2}-|\mathbf p_i|^2,
 \qquad a=1,2.
\end{equation*}
因此 \eqnrefs{eq:invariant-flux-scalar-r9} 根号中的量逐项化为
\begin{align*}
 \mathcal I_{12}^2
 &=\left(\frac{E_1E_2}{c^2}+p_i^2\right)^2
 -\left(\frac{E_1^2}{c^2}-p_i^2\right)
  \left(\frac{E_2^2}{c^2}-p_i^2\right),\\
 &=\frac{E_1^2E_2^2}{c^4}
 +\frac{2E_1E_2p_i^2}{c^2}+p_i^4
 -\frac{E_1^2E_2^2}{c^4}
 +\frac{(E_1^2+E_2^2)p_i^2}{c^2}-p_i^4,\\
 &=\frac{p_i^2}{c^2}(E_1+E_2)^2,
\end{align*}
其中 $p_i\equiv|\mathbf p_i|\ge0$。取正根得到
\begin{equation*}
 {
 \mathcal I_{12}
 =\frac{|\mathbf p_i|(E_1+E_2)}{c}
 =|\mathbf p_i|\sqrt s.}
 \end{equation*}
最后一步用了 $E_1+E_2=c\sqrt s$。于是
\begin{align*}
 v_{\rm M}
 &=|\mathbf p_i|c^2
 \left(\frac1{E_1}+\frac1{E_2}\right),\\
 &=|\mathbf v_1-\mathbf v_2|,
\end{align*}
因为 $\mathbf v_a=c^2\mathbf p_a/E_a$。所以 \eqnrefs{eq:moller-velocity-si-r9} 的协变表达确实回到熟悉的相对速度。
\end{derivation}

\begin{derivation}[从\eqnrefs{eq:delta-square-time-r15}到\eqnrefs{eq:2to2-cross-section-master-si-r9}]

取空间周期盒体积 $V$ 与观测时间 $T$。空间连续态满足
\begin{equation*}
 (2\pi\hbar)^3\delta^{(3)}(\mathbf0)=V,
 \qquad
 \sum_{\mathbf p}\longrightarrow
 V\int\frac{\dd^3 p}{(2\pi\hbar)^3}.
\end{equation*}
时间方向令 $\delta E\equiv E_f-E_i$，则
\begin{equation*}
 \int_{-T/2}^{T/2}\dd t\,\ee^{\ii\delta E t/\hbar}
 =T\,\operatorname{sinc}\!\left(\frac{\delta E T}{2\hbar}\right),
\end{equation*}
并由正文\eqnrefs{eq:delta-square-time-r15} 有
\begin{equation*}
 \lim_{T\to\infty}\frac1T
 \left|\int_{-T/2}^{T/2}\dd t\,\ee^{\ii\delta E t/\hbar}\right|^2
 =2\pi\hbar\delta(\delta E).
\end{equation*}
因此四维动量守恒 $\delta$ 函数平方在概率除以 $T$ 后留下一个能量守恒 $\delta$，其空间零点因子给出一个 $V$。结合两条初态的协变归一化，单位盒中单对入射粒子的二体末态跃迁率为正文\eqnrefs{eq:box-transition-rate-r15}，即
\begin{equation*}
 \dd\Gamma_{1+2\to2}
 =\frac{\hbar^2c^3}{4VE_1E_2}
 \overline{|\mathcal M|^2}\,\dd\mathfrak P_2.
\end{equation*}

单个归一化入射粒子在盒中的数密度是 $1/V$。正文\eqnrefs{eq:moller-velocity-si-r9} 给出相对速率，所以入射通量密度为
\begin{equation*}
 j_{12}=\frac{v_{\rm M}}{V}
 =\frac{c^3\mathcal I_{12}}{VE_1E_2}.
\end{equation*}
按正文中“跃迁率除以入射通量密度”的截面关系：
\begin{align*}
 \dd\sigma
 &=\frac{\dd\Gamma_{1+2\to2}}{j_{12}}\\
 &=\frac{\hbar^2}{4\mathcal I_{12}}
 \overline{|\mathcal M|^2}\,\dd\mathfrak P_2.
\end{align*}
$V$、$T$ 与参考系依赖的 $E_1E_2$ 在这一步全部消去。对 $2\to2$ 过程代入正文\eqnrefs{eq:phase2-final-si-r9}，并在质心系使用
$\mathcal I_{12}=|\mathbf p_i|\sqrt s$，得到
\begin{align*}
 \frac{\dd\sigma}{\dd\Omega}
 &=\frac{\hbar^2}{4|\mathbf p_i|\sqrt s}
 \overline{|\mathcal M|^2}
 \frac1{16\pi^2}\frac{|\mathbf p_f|}{\sqrt s}\\
 &=\frac{\hbar^2}{64\pi^2s}
 \frac{|\mathbf p_f|}{|\mathbf p_i|}
 \overline{|\mathcal M|^2},
\end{align*}
即正文\eqnrefs{eq:2to2-cross-section-master-si-r9}。
\end{derivation}

\begin{derivation}[从\eqnrefs{eq:delta-square-time-r15}到\eqnrefs{eq:master-decay-rate-si-r9}]
协变态归一化的单粒子初态在盒归一化后带因子 $(2P^0V)^{-1/2}$； 末态每条外腿给对应的质量壳测度。 LSZ 矩阵元中的四维 $\delta$ 在平方后，空间部分利用

\begin{equation*}
 (2\pi\hbar)^3\delta^{(3)}(\mathbf0)=V
\end{equation*}
消去母粒子的盒体积，而时间部分利用 \eqnrefs{eq:delta-square-time-r15} 把 $|\delta(E_f-E_i)|^2/T$ 化为单个能量 $\delta$。于是概率除以 $T$ 后只剩一个母粒子归一化因子 $1/(2P^0)$ 与 $n$ 体末态相空间。

为避免混淆能量量纲和动量量纲，先定义能量归一化截肢振幅
\begin{equation*}
 \mathcal A_E\equiv c\,\mathcal M.
\end{equation*}
若 $\mathcal M$ 在某个过程具有一个动量量纲，则 $\mathcal A_E$ 具有能量量纲。有限盒推导给
\begin{equation*}
 \dd\Gamma_{1\to n}
 =\frac{1}{2E_P\hbar}\,
 \overline{|\mathcal A_E|^2}\,\dd\mathfrak P_n,
 \qquad E_P=Mc^2.
 \end{equation*}
把 $\mathcal A_E=c\mathcal M$ 与 $E_P=Mc^2$ 代入，两个 $c^2$ 消去：
\begin{align*}
 \dd\Gamma_{1\to n}
 &=\frac{c^2}{2Mc^2\hbar}
 \overline{|\mathcal M|^2}\,\dd\mathfrak P_n,\\
 &={
 \frac{1}{2M\hbar}
 \overline{|\mathcal M|^2}\,\dd\mathfrak P_n.}
 \end{align*}
这里 $\dd\mathfrak P_n$ 使用能量四矢量 $\mathcal P^\mu=cp^\mu$ 的定义。以无量纲 Yukawa 耦合的费米子对衰变为例，振幅含有双线性因子 $\bar u v$；在上述外腿归一化下该双线性具有动量量纲，因此
\begin{equation*}
 \left[\frac{|\mathcal M|^2}{M\hbar}\right]
 =\mathrm{s^{-1}},
\end{equation*}
与衰变率的 SI 单位一致。这样 \eqnrefs{eq:master-decay-rate-si-r9} 与散射截面主式共享同一套有限盒、有限时间和外腿归一化，而不是独立记忆的另一条规则。
\end{derivation}



\begin{derivation}[从\eqnrefs{eq:phase3-definition-dp68}到\eqnrefs{eq:phase3-recursive-dp52}]
正文\eqnrefs{eq:phase3-definition-dp68} 是三体末态的原始 Lorentz 不变测度。内部不变质量 $\mathcal S_{12}$ 可以作为显式积分变量。引入 $\mathcal Q_{12}=\mathcal P_1+\mathcal P_2$，并插入

\begin{align*}
1
={}&\int\dd^4 \mathcal Q_{12}\,
\delta^{(4)}(\mathcal Q_{12}-\mathcal P_1-\mathcal P_2)\\
&\times\int\dd\mathcal S_{12}\,
\delta(\mathcal Q_{12}^2-\mathcal S_{12})\theta(\mathcal Q_{12}^0).
\end{align*}
正能质量壳恒等式
\begin{equation*}
\delta(\mathcal Q_{12}^2-\mathcal S_{12})\theta(\mathcal Q_{12}^0)
=\frac{1}{2E_{12}}\delta(\mathcal Q_{12}^0-E_{12})
\end{equation*}
把 $\mathcal Q_{12}$ 的四维积分化成标准单粒子质量壳测度。重新组合总四动量 $\delta$ 函数与三条末态测度后得到
\begin{align*}
\dd\mathfrak P_3
={}&\frac{\dd\mathcal S_{12}}{2\pi}
\left[\dd\mathfrak P_2(P;q_{12},3)\right]
\left[\dd\mathfrak P_2(q_{12};1,2)\right],
\end{align*}
这就是正文\eqnrefs{eq:phase3-recursive-dp52}。中间四动量 $q_{12}=p_1+p_2$ 仅用于把同一组三体末态按 $(12)+3$ 分组；其不变量 $q_{12}^2=\mathcal S_{12}$ 是三体相空间积分变量。
\end{derivation}

\begin{derivation}[从\eqnrefs{eq:phase3-recursive-dp52}到\eqnrefs{eq:phase3-dalitz-dp52}]
在母粒子静止系中，对正文\eqnrefs{eq:phase3-recursive-dp52} 的两个二体因子分别使用\eqnrefs{eq:phase2-final-si-r9}：

\begin{equation*}
\dd\mathfrak P_3
=\frac{\dd\mathcal S_{12}}{2\pi}
\frac{1}{16\pi^2}\frac{|\mathbf p_3|}{Mc}\dd\Omega_3
\frac{1}{16\pi^2}\frac{|\mathbf p_2^*|}{\sqrt{\mathcal S_{12}}/c}\dd\Omega_2^*.
\end{equation*}
若母粒子无自旋或已做自旋平均，整体方向积分给 $\int\dd\Omega_3=4\pi$。在 $q_{12}$ 静止系中以粒子 3 为极轴，方位角再给 $2\pi$，于是
\begin{equation*}
\dd\mathfrak P_3
=\frac{\dd\mathcal S_{12}\,\dd\cos\theta^*}{64\pi^3}
\frac{|\mathbf p_3|}{Mc}
\frac{|\mathbf p_2^*|}{\sqrt{\mathcal S_{12}}/c}.
\end{equation*}
正文\eqnrefs{eq:dalitz-angular-map-dp68} 给出
\begin{equation*}
\left|\frac{\partial\mathcal S_{23}}{\partial\cos\theta^*}\right|
=2c^2|\mathbf p_2^*||\mathbf p_3^*|.
\end{equation*}
因此把 $\cos\theta^*$ 换成 $\mathcal S_{23}$ 时，两个动量模恰好与 Jacobian 相消；再用 $\sqrt{\mathcal S_{12}}$ 和母粒子质量整理剩余因子，得到
\begin{equation*}
\dd\mathfrak P_3
=\frac{\dd\mathcal S_{12}\,\dd\mathcal S_{23}}{128\pi^3M^2c^4},
\end{equation*}
即正文\eqnrefs{eq:phase3-dalitz-dp52}。
\end{derivation}

\begin{derivation}[从\eqnrefs{eq:dalitz-angular-map-dp68}到\eqnrefs{eq:dalitz-boundary-dp52}]
在 $q_{12}$ 静止系中，两次二体运动学分别给

\begin{align*}
E_2^*
&=\frac{\mathcal S_{12}+m_2^2c^4-m_1^2c^4}{2\sqrt{\mathcal S_{12}}},\\
E_3^*
&=\frac{M^2c^4-\mathcal S_{12}-m_3^2c^4}{2\sqrt{\mathcal S_{12}}},\\
c|\mathbf p_2^*|
&=\frac{\sqrt{\Lambda_K(\mathcal S_{12},m_1^2c^4,m_2^2c^4)}}{2\sqrt{\mathcal S_{12}}},\\
c|\mathbf p_3^*|
&=\frac{\sqrt{\Lambda_K(M^2c^4,\mathcal S_{12},m_3^2c^4)}}{2\sqrt{\mathcal S_{12}}}.
\end{align*}
把这些量代入正文\eqnrefs{eq:dalitz-angular-map-dp68}。由于 $\mathcal S_{23}$ 对 $\cos\theta^*$ 是线性的，固定 $\mathcal S_{12}$ 时的最大值与最小值必定出现在 $\cos\theta^*=\mp1$。逐项代入便得到正文\eqnrefs{eq:dalitz-boundary-dp52}。

两个 Källén 根号还必须同时为实数，因此
\begin{equation*}
(m_1+m_2)^2c^4
\le\mathcal S_{12}
\le(M-m_3)^2c^4.
\end{equation*}
这给出 Dalitz 区域在 $\mathcal S_{12}$ 方向上的完整阈值范围。
\end{derivation}
